\documentclass[dvipdfmx]{beamer}
\usepackage[orientation=portrait,size=a0,scale=1.5]{beamerposter}%縦 : orientation=portrait, 横 : orientation=landscape
\usepackage{here, amsmath, latexsym, amssymb, bm, ascmac, mathtools, multicol, tcolorbox}
%カラーテーマの選択(省略可)
\usepackage{amsthm}
\usepackage{tikz}
\usepackage{pythonhighlight}
% \usetikzlibrary{plotmarks} % LATEX and plain TEX when using TikZ
% \usetikzlibrary{arrows.meta,positioning}
% \usetikzlibrary{datavisualization}
% \usetheme{Copenhagen}
% %フォントテーマの選択(省略可)
% \usefonttheme{professionalfonts}
% %フレーム内のテーマの選択(省略可)
% \useinnertheme{circles}
% %ナビゲーションバー非表示
% \setbeamertemplate{navigation symbols}{}
% %既定をゴシック体に
% \renewcommand{\kanjifamilydefault}{\gtdefault}
%itemize
% \setbeamertemplate{itemize item}{\small\raise0.5pt\hbox{$\bullet$}}
% \setbeamertemplate{itemize subitem}{\small\raise1.5pt\hbox{$\blacktriangleright$}}
% \setbeamertemplate{itemize subsubitem}{\small\raise1.5pt\hbox{$\bigstar$}}
% color
\newcommand{\red}[1]{\textcolor{red}{#1}}
\newcommand{\green}[1]{\textcolor{green!40!black}{#1}}
\newcommand{\blue}[1]{\textcolor{blue!80!black}{#1}}
\def \htheta{\hat{\theta}}
\def \cvT{{\cal C}_\mathrm{vT}}
\def \cBNH{{\cal C}_\mathrm{BNH}}
\def \cN{{\cal C}_\mathrm{N}}
\def \cHo{{\cal C}_\mathrm{H}}
\def \cP{{\cal C}_\mathrm{P}}
\def \cBSLD{{\cal C}_\mathrm{BSLD}}
\def \cRLD{{\cal C}_\mathrm{RLD}}
\def \cMI{{\cal C}_\mathrm{MI}}
\def \cGM{{\cal C}_\mathrm{BGM}}

\def \rhoB{S_{\rm B}}
\def \MBj{{D}_{\mathrm{B},j}}
\def \MBk{{D}_{\mathrm{B},k}}
\def \Mtheta{D(\theta)}
\def \Mthetaj{D_{j}(\theta)}
\def \MthetaT{D^\intercal(\theta)}
\def \mutwo{\mathsf{m}}
\def \muTwo{\mathsf{M}}
\def\sfA{\mathsf{A}}
\def\sfB{\mathsf{B}}
\def\sfc{\mathsf{c}}
\def\sfK{\mathsf{K}}
\def\sfM{\mathsf{M}}
\def\sfR{\mathsf{R}}
\def\sfJ{\mathsf{J}}
\def\sfV{\mathsf{V}}
\def\sfH{\mathsf{H}}
\def\sfS{\mathsf{S}}
\def\sfT{\mathsf{T}}
\def\sfZ{\mathsf{Z}}
\def\sfw{\mathsf{w}}
\def\sfW{\mathsf{W}}
\def\bbH{\mathbb{H}}
\def\bfS{\mathbf{S}}
\def\sfL{\mathsf{L}}
\def\bbU{\mathbb{U}}
\def\bbV{\mathbb{V}}
\def\bbX{\mathbb{X}}
\def \prior{\pi(\theta)}
\newcommand{\opTr}[1]{\mathrm{Tr}\left\{#1\right\}}
\newcommand{\sfTr}[1]{\mathrm{Tr}\left\{#1\right\}}
\newcommand{\bbTr}[1]{\mathrm{Tr}\left\{#1\right\}}
% \newtheorem{corollary}[theorem]{Corollary}
\newtheorem{prim}{prim}
\newtheorem{dual}{dual}

\def \bfSbar{\overline{\bfS}}
\def \Mbar{\overline{D}}
\def \Mbarj{\overline{D}_j}
\def \sfwbar{\overline{\sfw}}
\def \sfwtheta{\sfw(\theta)}% \newtheorem{corollary}[theorem]{Corollary}

\newcommand{\inner}[2]{\langle #1| #2\rangle}
\newcommand{\Inner}[2]{\langle #1, #2\rangle}

\newcommand{\ket}[1]{|#1\rangle}
\newcommand{\bra}[1]{\langle#1|}
%%%%% for draft only
%\usepackage{xcolor}
\usepackage{ulem}
%\newcommand{\red}[1]{{\color{red}{#1}}}
%\newcommand{\blue}[1]{\textcolor{blue}{#1}}
\newcommand{\rsout}[1]{\textcolor{red}{\sout{#1}}}
%\newcommand{\green}[1]{{\color{green}{#1}}}
\newcommand{\orange}[1]{\textcolor{orange}{#1}}

\setbeamertemplate{navigation symbols}{}
%head
\setbeamertemplate{headline}{
    \begin{center}
        \structure{
            \rule{\linewidth}{4mm}
            \vskip3ex
            \begin{columns}
                \begin{column}{0.70\textwidth}
                    \usebeamercolor{title in headline}{\textbf{\LARGE{\inserttitle}}}\\[2.5ex]
                    \usebeamercolor{author in headline}{\large{\insertauthor}}\\[1.2ex]
                    \usebeamercolor{institute in headline}{\normalsize{\insertinstitute}}   
                \end{column}
                \begin{column}{0.25\textwidth}
                    \includegraphics[width=0.48\linewidth]{UEC_logo2.png}\includegraphics[width=0.48\linewidth]{erato_logo.png}
                \end{column}    
           \end{columns}

            \vskip2ex
            \rule{\linewidth}{4mm}
            \vspace{-20pt}
        }
        \end{center}
}
%foot
\setbeamertemplate{footline}{
    \begin{center}
        \structure{
            \rule{0.98\linewidth}{2mm}
            \vskip4ex
            \begin{columns}[T]
                \begin{column}{0.05\paperwidth}
                \end{column}
                \begin{column}{0.95\paperwidth}
                    \usebeamercolor{conference in headline}{\small{\foot}}
                \end{column}
            \end{columns}
            \vskip3ex
        }
    \end{center}
}

\setbeamertemplate{theorems}[ams style]
\newtcolorbox{mybox}[1]
{
    title=#1, 
    toptitle=2mm, bottomtitle=2mm, 
    colframe=structure,boxrule=5pt,
    arc=30pt,
    coltitle=white, colbacktitle=structure,
    colback=white, fonttitle=\bfseries\Large,
    top=3mm, bottom=3mm, left=5mm, right=5mm, %内部余白調整
    enlarge top by=0mm, enlarge bottom by=0mm, %外部余白調整
}
\newtcolorbox{shortbox}{
    toptitle=2mm, bottomtitle=2mm, 
    colframe=structure,boxrule=5pt,
    arc=30pt,
    coltitle=white, colbacktitle=structure,
    colback=white, fonttitle=\bfseries\Large,
    top=3mm, bottom=3mm, left=1mm, right=1mm, %内部余白調整
    enlarge top by=0mm, enlarge bottom by=0mm, %外部余白調整
    }
\title{
A Comparison Between Bayesian Method and\\
~Linear Estimation in Qubit Tomography}
\author{Ke-Han Zhao†, Jun Suzuki†, E-mail: c2241016@gl.cc.uec.ac.jp}
\institute{† Graduate School of Informatics and Engineering, The University of Electro-Communications\\
}
\newcommand{\foot}{
    Acknowledgment:
    The work is partly supported by JSPS KAKENHI Grant Number JP24K14816, 
    and ERATO ``Super Quantum Entanglement" (Grant No. JPMJER2402) from JST.
    }

\begin{document}

\begin{frame}[fragile]
\begin{columns}[t]
    \begin{column}{0.5\linewidth}

        \begin{mybox}{Introduction}
            \begin{itemize}
                \item \textbf{Problem:} Quantum state estimate is unstable with limited sample size. 
                \item \textbf{Approach:} Use Bayesian framework with prior information.
                \item \textbf{Advantage of Bayesian statistics}:
                \begin{itemize}
                    \item[$\bullet$] Low variance in small-sample size
                    \item[$\bullet$] Posterior distribution produces uncertainty. 
                    \item[$\bullet$] Estimator is always physical ($|\theta| \leq 1$) with an appropriate prior. 
                \end{itemize}
            \end{itemize}

        \end{mybox}

        \begin{mybox}{Model and problem setting}
            \begin{itemize}
                \item \textbf{Parametric model}
                    \item[]  Single-qubit state parametrized by Bloch vector $\theta\in \mathbb{R}^3$:
                    \begin{equation}
                        \rho_\theta=\frac{1}{2}\left(I+\theta_1\sigma_x+\theta_2\sigma_y+\theta_3\sigma_z\right),\quad \sum_{i=1}^{3}\theta_i^2\leq1.
                    \end{equation}
                \item \textbf{Three measurement settings}  
                \item[] Measure in the X/Y/Z basis.
                \item[] Collected Measurement outcomes ($k_i$) for each basis obey the Bernoulli distribution:
                \begin{equation}
                    k_i \thicksim \mathrm{Binomial}(n_i,p_i), p_i=\frac{1}{2}(1+\theta_i) \quad \text{for } i=1,2,3
                \end{equation}
                \item \textbf{Problem}
                \item [] \hspace*{20mm} \blue{Bayesian methods}\hspace{20mm} vs. \hspace{25mm} \orange{Frequentist methods}
        \begin{columns}
            \begin{column}{0.5\linewidth}
                \begin{itemize}
                    \item [] \hspace*{-5mm}\blue{Expected a posteriori estimator (EAP)}
                    \item [] \hspace*{-5mm}\blue{Maximum a posteriori estimator (MAP)}
                \end{itemize}
            \end{column}
            \begin{column}{0.5\linewidth}
                \begin{itemize}

                    \item[] \hspace*{-3mm}\orange{Liner inversion estimator (LI)}
                    \item[] \hspace*{-3mm}\orange{Maximum likelihood estimator (MLE)}
                \end{itemize}
            \end{column}
        \end{columns}
        \ \\
            \item \textbf{Prior distribution in Bayesian framework}
            \item[] Parameter $\theta$ obeys a distribution:
            \begin{equation}
                    \theta=(\theta_1,\theta_2,\theta_3) = r \overrightarrow{u} 
            \end{equation}
            \item[$\bullet$] $r\thicksim \mathrm{Beta}(\alpha,\beta)$: purity ($\alpha,\beta>0$: hyperpameters)
            \item[$\bullet$] $\overrightarrow{u}\thicksim \mathrm{vMF}(\kappa,\mu)$: direction ($\kappa>0$, unit vector $\mu$: hyperparameters)
        \end{itemize}
        \end{mybox}
            
            \begin{mybox}{Simulation example 1}
                \begin{itemize}
                    \item \textbf{Simulation setting}
                    \item[] Two states: $|\theta|=0.99$ (close to pure), $|\theta|=0.6$ (more mixed).
                    \item[] $n$ measurements for each X/Y/Z basis with $n=10,50$.
                    \item \textbf{Setting hyperparameters}
                    \item [] As an example, $(\alpha, \beta)=(1,4)$, $\kappa=1$, and $\mu=(1,0,0)$, i.e., the state should close to $\frac{1}{2}(I+\sigma_x)$, with weak directional concentration and a broad purity distribution.
                \end{itemize}
        \begin{columns}
            \begin{column}{0.5\linewidth}
\begin{center}
                        \includegraphics[width=0.5\linewidth]{beta41.png}
    
\end{center}  
          \end{column}
            \begin{column}{0.5\linewidth}
\begin{center}
                            \includegraphics[width=0.5\linewidth]{Figure_vmf1.png}
\end{center}   
         \end{column}
        \end{columns}

                \begin{itemize}
                    \item  \textbf{Sampling via PyMC (Python library)}\\
                    \item[] It is easy to customize distribution with PyMC and do Sampling.
                \end{itemize}
                        \begin{columns}
            \begin{column}{0.7\linewidth}
                \begin{center}
                        \includegraphics[]{theta_mse_i1.png}
                \end{center}   
            \end{column}
            \begin{column}{0.3\linewidth}
                            \includegraphics[]{mapeap.png} 
            \end{column}
        \end{columns}

                \begin{columns}
                    \begin{column}{0.45\linewidth}
                        \begin{itemize}
                            \item EAP: stable estimate  
                            \item MAP: prior-sensitive 
                        \end{itemize}
                    \end{column}
                    \begin{column}{0.5\linewidth}
                        \begin{itemize}
                            \item MLE: unstable for mixed state
                            \item LI: often non-physical ($|\theta| > 1$)   
                        \end{itemize}
                    \end{column}
                \end{columns}
                \begin{shortbox}
                    \begin{center}
                        Bayesian method (EAP) outperforms maximum likelihood estimation for limited sample size.
                    \end{center}
                \end{shortbox}
            \end{mybox}

    

    \end{column}

    \begin{column}{0.5\linewidth}
        \begin{mybox}{Simulation example 2}
            \begin{itemize}
            \item \textbf{Less prior knowledge: Suppose we do not have enough information about $\theta$.}
                \item \textbf{Non-informative prior for the purity $r$ (Jeffereys prior)}
            \end{itemize}
        \begin{columns}
            \begin{column}{0.5\linewidth}
\begin{center}
                        \includegraphics[width=0.5\linewidth]{beta.5.5.png}
    
\end{center}  
          \end{column}
            \begin{column}{0.5\linewidth}
\begin{center}
                            \includegraphics[width=0.5\linewidth]{Figure_vmf3.png}
\end{center}   
         \end{column}
        \end{columns}
                \begin{center}
                    \includegraphics[width=0.7\linewidth]{theta_mse_i2.png}
                \end{center}
                \begin{itemize}
                \item[] EAP: Best under mixed state.
                \item[] MAP becomes unstable when the prior has unbounded density.
%                \begin{columns}
%                    \begin{column}{0.45\linewidth}
%                        \begin{itemize}
%                            \item EAP: Best under mixed state.
%                        \end{itemize}
%                    \end{column}
%                    \begin{column}{0.55\linewidth}
%                        \begin{itemize}
%                            \item MAP becomes unstable when the prior has unbounded density.  
%                        \end{itemize}
%                    \end{column}
%                \end{columns}
            \end{itemize}
            \begin{shortbox}
            \begin{center}
                \ EAP achieves lower variance under weakly informative prior.
            \end{center}
            \end{shortbox}
        \end{mybox}


        \begin{mybox}{Simulation example 3}

                \begin{itemize}
                \item \textbf{Prior misspecification: pior distribution is significantly different from the true one.}
                \item $(\alpha,\beta)=(1,4)$, $\kappa=4$, $\mu=(\red{-1},0,0)$, i.e., opposite direction.
                \end{itemize}
                        \begin{columns}
            \begin{column}{0.7\linewidth}
                \begin{center}
                        \includegraphics[]{theta_mse_i_oppsite.png}
                \end{center}   
            \end{column}
            \begin{column}{0.3\linewidth}
\begin{center}
                                \includegraphics[width=0.8\linewidth]{Figure_vmf2.png} 
    
\end{center}            
\end{column}
        \end{columns}
        \begin{itemize}
                \item[] Under severely misspecified prior, MAP degrades significantly, while EAP is less sensitive.
                \item A weak \red{secondary mode} of the posterior indicates prior-data mismatch.
                \begin{center}
                    \includegraphics[width=0.9\linewidth]{theta_trace_miss.png}
                \end{center}

        \end{itemize}
                \begin{shortbox}
                \begin{center}
                        Posterior distribution provides diagnostic insight.
                \end{center}                
                \end{shortbox} 
        %   \begin{shortbox}
        %     \begin{center}
        %         It is a projection-valued measure!
        %     \end{center}
        %   \end{shortbox}
        \end{mybox}
        \begin{mybox}{Conclusion and outlook}
            \begin{itemize}
                \item EAP yields stable estimates even for small sample size data.
                \item EAP is less sensitive to prior choice and achieves low variance.
                \item Posterior can provide uncertainty and suggest prior-data mismatch.
                \item[]\textbf{Outlook}
                \item Develop systematic approaches for hyperparameter selection, such as prior updating or other reliable methods.
                \item Extend to hierarchical Bayesian models to handle imbalanced or heterogeneous measurement settings.
        \end{itemize}

        \end{mybox}
            % \begin{shortbox}
            %     \small
            %     \begin{itemize}
            %         \item[{[1]}] R. D. Gill and S. Massar, State estimation for large ensembles, PRA, 61, 042312 (2000).
            %         \item[{[2]}] S.D. Personick, Application of quantum estimation theory to analog communication over quantum channels, IEEE T-IT,17, 240 (1971).
            %         \item[{[3]}] J. Rubio and J. Dunningham, Bayesian multiparameter quantum metrology with limited data, PRA, 101, 032114 (2020).
            %         \item[{[4]}] R. Demkowicz-Dobrza{\'n}ski, W. G{\'o}recki, and M. Gu{\c{t}}{\u{a}},Multi-parameter estimation beyond quantum Fisher information, JPhys A, 53, 363001 (2020).
            %         \item[{[5]}] K. Yamagata, Efficiency of quantum state tomography for qubits, IJQI, 9, 1167 (2011).
            %         \item [{[6]}]  J. Suzuki, Bayesian Nagaoka-Hayashi Bound for Multiparameter Quantum-State Estimation Problem, IEICE Trans, E107A, 510 (2024).
            %     \end{itemize}
                
            % \end{shortbox}
    \end{column}
\end{columns}
\end{frame}

\end{document}
